% ==============================================================================
% РАЗДЕЛ 2.2: Оценка переходного процесса и влияния различных факторов
% Партия B03-019 (DISS-005-B0813 .. DISS-005-B0861, 49 блоков)
% ==============================================================================

% DISS-005-B0813
\subsection{Оценка переходного процесса и влияния различных факторов}
\label{sec:ch02_transient_and_harmonics}

% DISS-005-B0814
Мы можем видеть (Рисунок 2.5), что для классических ТТ, при наличии апериодической составляющей в первичном токе, довольно длительное время сохраняются колебания (127~мс). Для катушек Роговского, данный недостаток практически отсутствует, что особенно ярко видно при построении годографа векторов тока (см. рис. ). Годограф отображает траекторию движения вектора, который используется для работы различных направленных защит, например, РНМ в устройствах БАВР.

% DISS-005-B0815
В случае применения КР, вектор из точки с нулевым током (точка 1, Рисунок 2.6) достигает установившегося значения (точка 2, Рисунок 2.6) за время, равное периоду промышленной частоты (20~мс). При этом траектория движения вектора, рассчитанного на основе этого сигнала, не зависит от угла КЗ.

% DISS-005-B0816
Хотя классические ТТ могут достигать близкого значения быстрее ($\varphi_{\text{КЗ}} = 0$)~-- за 15~мс (около точки 2, Рисунок 2.6)\commentnote{DISS-005-C0077}{Андрей Яблоков}{Вот это не сходится с рисунком вроде, все годографы кроме темно синего достигает точки 3 быстрее. Надо как то пояснить мысль, читатель не поймет}\commentnote{DISS-005-C0078}{Алексей Евдаков}{Переписал. Ибо суть в том, что у\textbackslash\textbackslash КР -- траектория стабильна и одинакова\textbackslash\textbackslash ТТ -- траектория изменяется от момента начала аварии}, но в этом случае вектор колеблется вокруг точки установившегося значения ещё длительное время, и входит в допустимый диапазон только через 127~мс от начала процесса (точка 2, Рисунок 2.5). При других углах КЗ время переходного процесса изменяется, и в случае отсутствия апериодической составляющей ($\varphi_{\text{КЗ}} = 90$) может быть аналогично с КР, то есть достигать установившегося значения (точка 2, рис. 8) за период промышленной частоты (20~мс).

% DISS-005-B0817
Таким образом, у классических ТТ, в зависимости от момента начала аварии ($\varphi_{\text{КЗ}}$), наблюдается непостоянство траектории движения вектора тока и времени переходного процесса. В отличие от них, применение КР позволяет избежать этой неоднозначности, обеспечивая более стабильное и предсказуемое поведение защит.

% DISS-005-B0818
% DISS-005-B0819
\begin{figure}[htbp]
  \centering
  \includegraphics[width=0.85\textwidth]{image60.png}
  \caption{Годограф вектора тока получаемого с классических ТТ при различных углах КЗ (CT\_0 -- CT\_90) с учётом корректировки по углу, с КР, а также усреднённого сигнала при максимальной доли аппериодики. Точка 1~-- начало процесса, точка 2~-- окончание процесса.}
  \label{fig:ch02_hodograph}
\end{figure}

% DISS-005-B0820
\paragraph{Влияние постоянной времени затухания ($\tau_0$).}

% DISS-005-B0821
Доля апериодической составляющей была задана максимальной ($\varphi_{\text{КЗ}} = 0$), а все иные параметры, участвующие в формуле (2.1)\commentnote{DISS-005-C0079}{Андрей Яблоков}{Надо описать где то, что это за иные параметры исходной модели}\commentnote{DISS-005-C0080}{Алексей Евдаков}{Пока не придумал, как лучше. Да и\dots\ других параметров особо нет. Это скорее совсем базовые, например частота сети, длина процесса и т.д.}\commentnote{DISS-005-C0081}{Андрей Яблоков}{Может сослаться на формулу: все иные параметры тока, участвующие в формуле 1, \dots.}\commentnote{DISS-005-C0082}{Алексей Евдаков}{Да, спасибо за идею.}, оставлены исходными для данной модели. Как можем видеть, на всём промежутке значений $\tau_0$: от 1~мс до 700~мс, что охватывает все встречаемые на практике значения параметра (от 1~мс до 540~мс согласно [24]), использование КР даёт либо лучшее время переходного процесса, либо соизмеримое (Рисунок 2.7~а,~б). Причём, при постоянных времени более 40~мс, время переходного процесса определяется исключительно математическими факторами и особенностями БПФ и равно 20~мс. Классические же ТТ могут достигать времени переходного процесса равного порядка 240~мс. При этом, стоит не забывать, что трансформаторы тока с ферромагнитными сердечниками часто успевают насытиться за это время, и их показания окажутся ещё хуже, чем представленные здесь для идеального сигнала.

% DISS-005-B0822
\paragraph{Влияние начального угла КЗ.}

% DISS-005-B0823
Переменная $\tau_0$ была задана равной 50~мс, все остальные параметры оставлены стандартными. Можем видеть (Рисунок 2.7~в,г), что чем больше доля апериодической составляющей (чем ближе угол КЗ к 0), тем это искажение сильнее у классических ТТ. Для КР же, при данной $\tau_0$, время переходного процесса определяется исключительно математическими факторами, как и было показано ранее.

% DISS-005-B0824
\paragraph{Влияние частоты дискретизации.}

% DISS-005-B0825
По результатам моделирования получено, что количество выборок на период не оказывает существенного влияния на точность измерения при стандартных значениях (20, 24, 32, 48, 96 и более).

% DISS-005-B0826
% DISS-005-B0827
% DISS-005-B0828
% DISS-005-B0829
% DISS-005-B0830
% DISS-005-B0831
% DISS-005-B0832
% DISS-005-B0833
% DISS-005-B0834
% DISS-005-B0835
% DISS-005-B0836
\begin{figure}[htbp]
  \centering
  \begin{tabular}{cc}
    \includegraphics[width=0.46\textwidth]{image61.png} & \includegraphics[width=0.46\textwidth]{image62.png} \\
    а) & б) \\
    \includegraphics[width=0.46\textwidth]{image63.png} & \includegraphics[width=0.46\textwidth]{image64.png} \\
    в) & г)
  \end{tabular}
  \caption{Время переходного процесса:\\
  -- при различном времени затухания аппериодической составляющей: а) определяемое по модулю, б) определяемое по углу;\\
  -- при различном начальном угле КЗ: в) определяемое по модулю, г) определяемое по углу}
  \label{fig:ch02_transient_time_factors}
\end{figure}

% DISS-005-B0837
\subsection{Влияние высших гармоник}
\label{sec:ch02_higher_harmonics}

% DISS-005-B0838
Физическая природа катушек Роговского (взятие производной) приводит к увеличению сигналов высших частот, поэтому является важной оценка влияния высших гармоник в сигнале на расчёт амплитуды и угла первой гармоники.

% DISS-005-B0839
Исследование данного фактора осуществлялось путём добавления в сигнал 2, 3, и 10 гармоник, базовые амплитуды которых заданы 0,3, 0,2 и 0,1 о.е. соответственно. Вторая и третья гармоники описывают характерные для реальной сети искажения, вызванные насыщение силовых устройств (трансформаторов и двигателей), а десятая гармоника используется для моделирования шумовой составляющей. Стартовый угол КЗ для всех гармоник задавался одинаковым и равным 0 градусов, это было выбрано для максимизации\commentnote{DISS-005-C0083}{Valery Valery}{Это слово хорошо переведется на английский, но оно просто ужасно))} влияния апериодической составляющей. Формулы используемых сигналов приведены ниже (2.5):

% DISS-005-B0840
% DISS-005-B0841
% DISS-005-B0842
{\small
\begin{equation}
\label{eq:ch02_harmonics_signals}
\begin{cases}
i(t) = I_{\text{кз}} \cdot \sqrt{2} \cdot \left( (1 + h_2 + h_3 + h_{10}) \cdot \cos(\varphi_{\text{КЗ}}) \cdot e^{-\frac{t}{\tau_0}} - f(t) \right) \\
f(t) = \cos(\omega t + \varphi_{\text{КЗ}}) - h_2 \cdot \cos(2\omega t + \varphi_{\text{КЗ}}) - h_3 \cdot \cos(3\omega t + \varphi_{\text{КЗ}}) - h_{10} \cdot \cos(10\omega t + \varphi_{\text{КЗ}}) \\
i'(t) = I_{\text{кз}} \cdot \sqrt{2} \cdot \left( -\frac{\cos(\varphi_{\text{КЗ}})}{\tau_0} \cdot (1 + h_2 + h_3 + h_{10}) \cdot e^{-\frac{t}{\tau_0}} + f'(t) \right) \\
\begin{aligned}
f'(t) = \omega t \cdot \sin(\omega t + \varphi_{\text{КЗ}}) &+ h_2 \cdot 2\omega t \cdot \sin(2\omega t + \varphi_{\text{КЗ}}) +{} \\
&+ h_3 \cdot 3 \cdot 2\omega t \cdot \sin(3\omega t + \varphi_{\text{КЗ}}) + h_{10} \cdot 10 \cdot 2\omega t \cdot \sin(10\omega t + \varphi_{\text{КЗ}})
\end{aligned}
\end{cases}
\end{equation}
}

% DISS-005-B0843
где $i(t)$ и $i'(t)$~-- мгновенные значения получаемых сигналов исходного сигнала и его производной, $I_{\text{кз}}$~-- амплитуда тока КЗ, $\varphi_{\text{КЗ}}$~-- угол начала КЗ, определяющий долю аппериодической составляющей, $\tau_0$~-- постоянная времени затухания, $\omega = 2 \cdot \pi \cdot f$~-- угловая частота, $t$~-- время переходного процесса, $f(t)$ и $f'(t)$~-- периодические составляющие сигналов.

% DISS-005-B0844

% DISS-005-B0845
% DISS-005-B0846
\begin{figure}[htbp]
  \centering
  \includegraphics[width=0.88\textwidth]{image65.png}
  \caption{Исследуемые сигналы с ТТ без учёта насыщения и КР с наличием 2, 3, и 10 гармоник в сигнале; и рассчитанные амплитуды первой гармоники.}
  \label{fig:ch02_harmonics_currents}
\end{figure}

% DISS-005-B0847
На рисунке (Рисунок 2.8) показаны исходные сигналы, и на их основе рассчитанные амплитуды первых гармоник. Как можем видеть, наличие высших гармонических составляющих для КР вносит дополнительное искажение только на первом периоде расчёта, далее амплитуда остаётся стабильной. Данное поведение и ожидалось, так как описанные неточности измерения обусловлены апериодической составляющей, которая лишь немного изменилась по амплитуде.

% DISS-005-B0848
\subsection{Ускорение работы алгоритмов на начальном участке и применение КР для алгоритмов на мгновенных значениях}
\label{sec:ch02_instantaneous_values}

% DISS-005-B0849
Как можно было заметить ранее, на начальном участке значение амплитуды первой гармоники, рассчитанное на основе сигнала с катушки Роговского, опережает аналогичное значение с идеальных классических ТТ (Рисунок 2.9). Данный фактор может быть полезен при функционировании быстродействующих устройств РЗА и БАВР в первые моменты развития аварии.

% DISS-005-B0850
% DISS-005-B0851
% DISS-005-B0852
% DISS-005-B0853
% DISS-005-B0854
\begin{figure}[htbp]
  \centering
  \begin{tabular}{cc}
    \includegraphics[width=0.48\textwidth]{image66.png} & \includegraphics[width=0.48\textwidth]{image67.png} \\
    а) & б)
  \end{tabular}
  \caption{Переходный процесс, амплитуды первой гармоники: а) при максимальной доле аппериодической составляющей и постоянной времени 50~мс, б) при отсутствии аппериодической составляющей}
  \label{fig:ch02_initial_transient}
\end{figure}

% DISS-005-B0855
Алгоритмы получения амплитуды сигнала, основанные на мгновенных значениях [22]\commentnote{DISS-005-C0084}{Алексей Евдаков}{Добавил одну работу, по-хорошему бы ещё добавить методичку, по которой делается лабораторная, о которой говорил Андрей.}\commentnote{DISS-005-C0085}{Евдаков Алексей}{Эту статью надо обновить / вставить.}, могут быть использованы для ускорения функционирования устройств. \commentnote{DISS-005-C0086}{Андрей Яблоков}{Не надо начинать предложения с также}Если для классических ТТ у данных подходов возникают существенные погрешности при апериодической составляющей, то для катушек Роговского данный эффект практически отсутствует ввиду её подавления физической природой датчика.

% DISS-005-B0856
Таким образом, применение катушек Роговского в устройствах релейной защиты и автоматики обеспечивает:
\begin{itemize}[noitemsep,topsep=2pt]
% DISS-005-B0857
\item Устойчивость работы защит за счет исключения насыщения, характерного для ТТ.
% DISS-005-B0858
\item Снижение влияния апериодической составляющей на расчет первой гармоники.
% DISS-005-B0859
\item Сокращение времени расчётного переходного процесса, что критично для быстродействующих защитных устройств.
% DISS-005-B0860
\item Экономическую эффективность и простоту монтажа, особенно при модернизации существующих подстанций.
\end{itemize}

% DISS-005-B0861
При этом совместное использование сигналов с классических ТТ и КР является интересной идей для поиска баланса между данными типами датчиков. В дальнейшем предлагается использовать нелинейную зависимость, при которой на малых токах приоритет отдается сигналу с ТТ для повышения точности за счет исключения шумов, а при больших - сигналу с КР для компенсации насыщения сердечника ТТ. Такая концепция может быть применима, например, на цифровых, так как приём сигнала по SV потоку не требует прокладки дополнительных кабелей. Разные типы датчиков могут присутствовать в комбинированных измерительных приборах, например: ТТ используется для измерений, а КР~-- для целей РЗА.
